\setcounter{section}{6}

  \zwischenueberschrift{Turunan Kovarian}

Untuk suatu hipermuka terdiferensialkan

\mavergleichskette
{\vergleichskette
{ Y
}
{ \subseteq }{ \R^n
}
{  }{
}
{    }{
}
{   }{
}
}
{}{}{}
kita ingin mendefinisikan konsep-konsep turunan yang mudah didefinisikan di ruang sekitar $\R^n$, tetapi memiliki kaitan alami dengan $Y$. Konsep ini disebut turunan kovarian; alat bantu terpentingnya adalah proyeksi ortogonal
\maabbdisp {\pi_P} {\R^n } { T_PY
} {}
di suatu titik

\mavergleichskette
{\vergleichskette
{ P
}
{ \in }{ Y
}
{  }{
}
{    }{
}
{   }{
}
}
{}{}{.}

\emph{Catatan edisi: Relasi pembuka diperbaiki dari $Y=\R^n$ menjadi $Y\subseteq\R^n$; sebuah hipermuka tidak sama dengan seluruh ruang sekitarnya.}

\inputdefinition
{}
{

Misalkan

\mavergleichskette
{\vergleichskette
{ Y
}
{ \subseteq }{ W
}
{ \subseteq }{ \R^n
}
{    }{
}
{   }{
}
}
{}{}{,}
$W$
\definitionsverweis {terbuka}{}{,}
dan $Y$ suatu
\definitionsverweis {hipermuka terdiferensialkan}{}{,}
misalkan

\mavergleichskette
{\vergleichskette
{ P
}
{ \in }{ Y
}
{  }{
}
{    }{
}
{   }{
}
}
{}{}{}
suatu titik dan

\mavergleichskette
{\vergleichskette
{ v
}
{ \in }{ T_PY
}
{  }{
}
{    }{
}
{   }{
}
}
{}{}{}
suatu
\definitionsverweis {vektor singgung}{}{.}
Misalkan
\maabbdisp {F} { U} { \R^n
} {}
suatu medan vektor terdiferensialkan yang didefinisikan pada lingkungan terbuka

\mavergleichskette
{\vergleichskette
{ P
}
{ \in }{ U
}
{ \subseteq }{ \R^n
}
{    }{
}
{   }{
}
}
{}{}{}
dan pada $Y \cap U$ tangensial terhadap $Y$. Maka

\mavergleichskettedisp
{\vergleichskette
{ \nabla_v F
}
{ \defeq} { \pi_P   \left(  \left(  D F   \right)_{ P } \left(  v  \right)  \right)
}
{ } {
}
{ } {
}
{ } {
}
}
{}{}{,}
dengan
\maabbdisp {\pi_P} {\R^n } { T_PY
} {}
menyatakan
\definitionsverweis {proyeksi ortogonal}{}{,}
disebut
\definitionswort {turunan kovarian}{}
dari $F$ dalam arah $v$.

}

Secara keseluruhan, situasinya adalah

\mathdisp {T_PY \longrightarrow \R^n \stackrel{  \left(D F \right)_{ P} }{\longrightarrow} \R^n\stackrel{ \pi_P }{\longrightarrow} T_PY} {  }

dan $\nabla_v F$ merupakan hasil ketika vektor singgung

\mavergleichskette
{\vergleichskette
{ v
}
{ \in }{ T_PY
}
{  }{
}
{    }{
}
{   }{
}
}
{}{}{}
dimasukkan pada awal rangkaian tersebut. Jika $N$ menyatakan suatu medan normal satuan, maka turunan kovarian sama dengan

\mavergleichskettedisp
{\vergleichskette
{ \nabla_v F
}
{ =} { \left(  D F   \right)_{ P } \left(  v  \right) -  \left\langle  \left(  D F   \right)_{ P } \left(  v  \right)  ,  N(P)  \right\rangle N(P)
}
{ } {
}
{ } {
}
{ } {
}
}
{}{}{,}
karena dengan cara inilah komplemen ortogonal dihitung.

\inputdefinition
{}
{

Misalkan

\mavergleichskette
{\vergleichskette
{ Y
}
{ \subseteq }{ W
}
{ \subseteq }{ \R^n
}
{    }{
}
{   }{
}
}
{}{}{,}
$W$
\definitionsverweis {terbuka}{}{,}
dan $Y$ suatu
\definitionsverweis {hipermuka terdiferensialkan}{}{}
dan misalkan
\maabbdisp {G} { Y } { TY
} {}
suatu medan vektor tangensial pada $Y$. Misalkan
\maabbdisp {F} { U} { \R^n
} {}
suatu medan vektor terdiferensialkan yang didefinisikan pada lingkungan terbuka

\mavergleichskette
{\vergleichskette
{ P
}
{ \in }{ U
}
{ \subseteq }{ \R^n
}
{    }{
}
{   }{
}
}
{}{}{}
dan pada $Y \cap U$ tangensial terhadap $Y$. Maka medan vektor tangensial
\maabbdisp {\nabla_G F} { Y \cap U } { T(Y \cap U)
} {,}
yang didefinisikan oleh

\mavergleichskettedisp
{\vergleichskette
{ \left(  \nabla_G F  \right) (P)
}
{ \defeq} { \nabla_{G(P)} F }
{ } {
}
{ } {
}
{ } {
}
}
{}{}{}
disebut
\definitionswort {turunan kovarian}{}
dari $F$ dalam arah $G$.

}

__NOEDITSECTION__

\inputfaktbeweis
{Hyperfläche/Tangentiales Vektorfeld/Bezüglich Tangentenvektor/Kovariante Ableitung/Eigenschaften/Fakt}
{Lemma}
{}
{

\faktsituation {Misalkan

\mavergleichskette
{\vergleichskette
{ Y
}
{ \subseteq }{ W
}
{ \subseteq }{\R^n
}
{    }{
}
{   }{
}
}
{}{}{,}
$W$
\definitionsverweis {terbuka}{}{,}
dan $Y$ suatu
\definitionsverweis {hipermuka terdiferensialkan}{}{,}
dan misalkan

\mavergleichskette
{\vergleichskette
{ P
}
{ \in }{ Y
}
{  }{
}
{    }{
}
{   }{
}
}
{}{}{}
suatu titik.}
\faktuebergang {Maka pernyataan-pernyataan berikut berlaku.}
\faktfolgerung {\aufzaehlungdrei{Untuk suatu
\definitionsverweis {medan vektor tangensial}{}{}
terdiferensialkan $F$ yang tetap, pemetaan
\maabbeledisp {} {T_PY} { T_PY
} { v } { \nabla_v F
} {,}
bersifat
\definitionsverweis {linear}{}{.}
}{Untuk suatu vektor singgung tetap

\mavergleichskette
{\vergleichskette
{ v
}
{ \in }{ T_PY
}
{  }{
}
{    }{
}
{   }{
}
}
{}{}{}
pemetaan
\maabbdisp {} { F } { \nabla_v F
} {,}
yang memetakan sebuah medan vektor terdiferensialkan ke turunan kovariannya sepanjang $v$, bersifat linear.
}{Untuk suatu fungsi terdiferensialkan
\maabb {g} { U } {\R
} {}
dan suatu medan vektor tangensial terdiferensialkan $F$, berlaku

\mavergleichskettedisp
{\vergleichskette
{ \nabla_v (gF)
}
{ =} { g (P) \nabla_v F + \left(  D_{v}  g   \right) \left(   P   \right) F(P)
}
{ } {
}
{ } {
}
{ } {
}
}
{}{}{.}
\emph{Catatan edisi: Evaluasi berlebih $(P)$ setelah $\nabla_vF$ dihapus karena $v\in T_PY$ sudah menentukan titiknya.}
}}
\faktzusatz {}
\faktzusatz {}

}
{

\aufzaehlungdrei{Hal ini langsung mengikuti representasi

\mavergleichskettedisp
{\vergleichskette
{ \nabla_v F
}
{ =} { \left(  D F   \right)_{ P } \left(  v  \right) - \left\langle  \left(  D F   \right)_{ P } \left(  v  \right)  ,  N(P)  \right\rangle N(P)
}
{ } {
}
{ } {
}
{ } {
}
}
{}{}{,}
karena kedua suku tersebut linear dalam $v$.
}{Hal ini juga mengikuti representasi

\mavergleichskettedisp
{\vergleichskette
{ \nabla_v F
}
{ =} { \left(  D F   \right)_{ P } \left(  v  \right) - \left\langle  \left(  D F   \right)_{ P } \left(  v  \right)  ,  N(P)  \right\rangle N(P)
}
{ } {
}
{ } {
}
{ } {
}
}
{}{}{,}
karena untuk vektor tetap $v$, menurut
[[Richtungsableitung/R/Elementare Eigenschaften/Fakt|Lemma 43.6 (Analisis (Osnabrück 2021--2023))]],
kedua suku bergantung secara linear pada $F$.
}{Menurut
[[Totale Differenzierbarkeit/K/Produktregel/Fakt|Lemma 45.10 (Analisis (Osnabrück 2021--2023))]],
berlaku

\mavergleichskettedisp
{\vergleichskette
{ \left(D( g F)\right)_{ P }
}
{ =} { g(P) \cdot \left(D F \right)_{ P } + \left(D g \right)_{ P } \cdot F (P)
}
{ } {
}
{ } {
}
{ } {
}
}
{}{}{.}
Oleh karena itu,

\mavergleichskettealignhandlinks
{\vergleichskettealignhandlinks
{ \nabla_v (gF)
}
{ =} { \left(  D(gF)  \right)_{ P } \left(  v  \right) - \left\langle  \left(  D(gF)  \right)_{ P } \left(  v  \right)  ,  N(P)  \right\rangle N(P)
}
{ =} { \begin{aligned}[t]
&g(P) \cdot \left(  D F   \right)_{ P } \left(  v  \right) + \left(  D g   \right)_{ P } \left(  v  \right) \cdot F (P) \\
&\quad - \left\langle  g(P) \cdot \left(  D F   \right)_{ P } \left(  v  \right) + \left(  D g   \right)_{ P } \left(  v  \right) \cdot F (P)  ,  N(P)  \right\rangle N(P)
\end{aligned}
}
{ =} { \begin{aligned}[t]
&g(P) \cdot \left(  D F   \right)_{ P } \left(  v  \right) - \left\langle  g(P) \cdot \left(  D F   \right)_{ P } \left(  v  \right)  ,  N(P)  \right\rangle N(P) \\
&\quad + \left(  D g   \right)_{ P } \left(  v  \right) \cdot F (P) - \left\langle  \left(  D g   \right)_{ P } \left(  v  \right) \cdot F (P)  ,  N(P)  \right\rangle N(P)
\end{aligned}
}
{ =} { g (P) \nabla_v F + \left(  D_{v}  g   \right) \left(   P   \right) \cdot F(P)
}
}
{}
{}{,}
karena medan vektor $F$ tegak lurus terhadap $N$.
}

}

__NOEDITSECTION__

\inputfaktbeweis
{Hyperfläche/Tangentiales Vektorfeld/Bezüglich tangentialem Vektorfeld/Kovariante Ableitung/Eigenschaften/Fakt}
{Lemma}
{}
{

\faktsituation {Misalkan

\mavergleichskette
{\vergleichskette
{ Y
}
{ \subseteq }{ W
}
{ \subseteq }{\R^n
}
{    }{
}
{   }{
}
}
{}{}{,}
$W$
\definitionsverweis {terbuka}{}{,}
dan $Y$ suatu
\definitionsverweis {hipermuka terdiferensialkan}{}{.}}
\faktuebergang {Maka pernyataan-pernyataan berikut berlaku bagi medan-medan vektor tangensial pada $Y$.}
\faktfolgerung {\aufzaehlungdrei{Untuk suatu medan vektor tangensial terdiferensialkan $F$ yang tetap, pemetaan
\mathl{G \mapsto \nabla_G F}{}
bersifat
\definitionsverweis {linear}{}{}
dalam $G$. Selanjutnya, untuk suatu fungsi kontinu
\maabb {g} { Y } {\R
} {}

\mavergleichskettedisp
{\vergleichskette
{ \nabla_{gG} F
}
{ =} { g \nabla_G F
}
{ } {
}
{ } {
}
{ } {
}
}
{}{}{.}
}{Untuk suatu medan vektor tangensial $G$ yang tetap, pemetaan
\mathl{F \mapsto \nabla_G F}{,} yang memetakan sebuah medan vektor tangensial terdiferensialkan ke turunan kovariannya sepanjang $G$, bersifat linear.
}{Untuk suatu fungsi terdiferensialkan
\maabb {g} { U } {\R
} {}
dan suatu medan vektor tangensial terdiferensialkan $F$, berlaku

\mavergleichskettedisp
{\vergleichskette
{ (\nabla_G (gF))(P)
}
{ =} { g (P) ( \nabla_G F) (P) + (\nabla_G g )(P) F(P)
}
{ } {
}
{ } {
}
{ } {
}
}
{}{}{.}
\emph{Catatan edisi: Ruas kiri dievaluasi di $P$ agar identitas ini konsisten dengan ruas kanan yang juga dievaluasi di $P$.}
}}
\faktzusatz {}
\faktzusatz {}

}
{

Semua pernyataan mengikuti dari
[[Hyperfläche/Tangentiales Vektorfeld/Bezüglich Tangentenvektor/Kovariante Ableitung/Eigenschaften/Fakt|Lemma 6.3]].

}

Pada persamaan terakhir,

\mavergleichskettedisp
{\vergleichskette
{ \left(  \nabla_G g  \right) (P)
}
{ =} { \left(  D_{ G(P)}  g   \right) \left(   P   \right)
}
{ } {
}
{ } {
}
{ } {
}
}
{}{}{,}
yang berarti turunan berarah fungsi $g$ diambil dalam arah $G(P)$.

\inputdefinition
{}
{

Misalkan

\mavergleichskette
{\vergleichskette
{ Y
}
{ \subseteq }{ W
}
{ \subseteq }{\R^n
}
{    }{
}
{   }{
}
}
{}{}{,} $W$
\definitionsverweis {terbuka}{}{,}
dan $Y$ suatu
\definitionsverweis {hipermuka terdiferensialkan}{}{,}
misalkan
\maabbdisp {\gamma} {I} {Y
} {}
suatu
\definitionsverweis {kurva terdiferensialkan}{}{}
dan misalkan
\maabbdisp {F} { I } { \R^n
} {}
suatu medan vektor terdiferensialkan sepanjang $\gamma$ yang tangensial terhadap $Y$. Maka

\mavergleichskettedisp
{\vergleichskette
{ (\nabla_\gamma F) (t)
}
{ \defeq} { \pi_{\gamma(t)} F' (t)
}
{ } {
}
{ } {
}
{ } {
}
}
{}{}{,}
disebut
\definitionswort {turunan kovarian}{}
dari $F$ sepanjang $\gamma$.

\emph{Catatan edisi: Domain $\gamma$ dan $F$ diseragamkan menjadi interval $I$, dan $F$ dinyatakan sebagai medan sepanjang $\gamma$; sumber memakai $[a,b]$ dan kemudian interval $I$ yang belum didefinisikan.}

}

Medan vektor $F$ merupakan kurva bernilai vektor di $\R^n$ dengan

\mavergleichskette
{\vergleichskette
{ F(t)
}
{ \in }{ T_{\gamma(t)}Y
}
{  }{
}
{    }{
}
{   }{
}
}
{}{}{.}
Turunan kovariannya diperoleh dengan menurunkan kurva ini lalu mengambil proyeksi ortogonal hasilnya ke ruang singgung. Kadang-kadang ditulis pula $\nabla_{\gamma'} F$ atau $\nabla_{\partial} F$\zusatzfussnote {Untuk
[[Hyperfläche/Tangentiales Vektorfeld/Bezüglich tangentialem Vektorfeld/Kovariante Ableitung/Definition|Definisi 6.2]]
dan
[[Hyperfläche/Tangentiales Vektorfeld/Längs einer Kurve/Kovariante Ableitung/Definition|Definisi 6.5]]
terdapat generalisasi berikut. Misalkan $Y$ adalah hipermuka terdiferensialkan
\zusatzklammer {atau suatu manifold diferensiabel dengan koneksi pada bundel tangen} {} {,}
misalkan $Z$ suatu manifold diferensiabel, misalkan
\maabb {\varphi} { Z } { Y
} {}
suatu pemetaan terdiferensialkan, misalkan
\maabbdisp {F} { Z } { \R^n
} {}
suatu pemetaan terdiferensialkan dengan $F(Q)\in T_{\varphi(Q)}Y$ untuk setiap $Q\in Z$, yaitu suatu medan vektor tangensial sepanjang $\varphi$, dan $G$ suatu medan vektor pada $Z$. Maka
\zusatzklammer {untuk

\mavergleichskette
{\vergleichskette
{ Q
}
{ \in }{ Z
}
{  }{
}
{    }{
}
{   }{
}
}
{}{}{}
dan

\mavergleichskette
{\vergleichskette
{ P
}
{ = }{ \varphi(Q)
}
{  }{
}
{    }{
}
{   }{
}
}
{}{}{}} {} {}
turunan kovarian $F$ terhadap $G$
\zusatzklammer {sepanjang $\varphi$} {} {}
didefinisikan oleh

\mavergleichskettedisp
{\vergleichskette
{ ( \nabla_G F )(Q)
}
{ =} { \pi_{ P  } \left(D F \right)_{ Q}  (G(Q))
}
{ } {
}
{ } {
}
{ } {
}
}
{}{}{}
dan $\nabla_GF$ kembali merupakan suatu pemetaan
\maabb {} { Z } { TY
}{.}
\emph{Catatan edisi: Kodomain $F$ diperbaiki menjadi $\R^n$ dengan syarat nilai tangensial, sehingga diferensial $DF$ yang diproyeksikan pada rumus di atas bertipe benar.}
Dalam definisi pertama,

\mavergleichskette
{\vergleichskette
{ Z
}
{ = }{ Y
}
{  }{
}
{    }{
}
{   }{
}
}
{}{}{}
dan $\varphi$ adalah identitas; dalam definisi kedua,

\mavergleichskette
{\vergleichskette
{ Z
}
{ = }{ I
}
{  }{
}
{    }{
}
{   }{
}
}
{}{}{}
adalah suatu interval,

\mavergleichskette
{\vergleichskette
{ \varphi
}
{ = }{ \gamma
}
{  }{
}
{    }{
}
{   }{
}
}
{}{}{}
adalah lintasan tersebut, dan

\mavergleichskette
{\vergleichskette
{ G
}
{ = }{ \partial
}
{  }{
}
{    }{
}
{   }{
}
}
{}{}{}
adalah medan vektor konstan pada interval $I$ dengan arah $1$.} {} {.}

  \zwischenueberschrift{Medan Vektor Paralel}

\inputdefinition
{}
{

Misalkan

\mavergleichskette
{\vergleichskette
{ Y
}
{ \subseteq }{ W
}
{ \subseteq }{ \R^n
}
{    }{
}
{   }{
}
}
{}{}{,}
$W$
\definitionsverweis {terbuka}{}{,}
dan $Y$ suatu
\definitionsverweis {hipermuka terdiferensialkan}{}{}
dan misalkan
\maabb {\gamma} { I } { Y
} {}
suatu
\definitionsverweis {kurva terdiferensialkan}{}{.}
Suatu medan vektor tangensial terdiferensialkan yang didefinisikan sepanjang $\gamma$,
\maabb {F} { I } {\R^n
} {}
disebut
\definitionswort {paralel sepanjang}{}
$\gamma$ jika

\mavergleichskettedisp
{\vergleichskette
{ (\nabla_{\gamma} F)(t)
}
{ =} { 0
}
{ } {
}
{ } {
}
{ } {
}
}
{}{}{}
untuk semua

\mavergleichskette
{\vergleichskette
{ t
}
{ \in }{ I
}
{  }{
}
{    }{
}
{   }{
}
}
{}{}{.}

}

Suatu medan vektor paralel sepanjang $\gamma$ tepat ketika
\mathl{F'(t)}{} selalu tegak lurus terhadap ruang singgung
\mathl{T_{\gamma(t)}Y}{}, yakni ketika turunannya berada pada garis normal.

__NOEDITSECTION__
\inputfaktbeweis
{Hyperfläche/Kurve/Paralleles Vektorfeld/Eigenschaften/Fakt}
{Lemma}
{}
{

\faktsituation {Misalkan

\mavergleichskette
{\vergleichskette
{ Y
}
{ \subseteq }{ W
}
{ \subseteq }{ \R^n
}
{    }{
}
{   }{
}
}
{}{}{,}
$W$
\definitionsverweis {terbuka}{}{,}
dan $Y$ suatu
\definitionsverweis {hipermuka terdiferensialkan}{}{}
dan misalkan
\maabb {\gamma} { I } { Y
} {}
suatu
\definitionsverweis {kurva terdiferensialkan}{}{.}}
\faktuebergang {Maka pernyataan-pernyataan berikut berlaku.}
\faktfolgerung {\aufzaehlungzwei {Jika $F$ adalah medan vektor
\definitionsverweis {paralel}{}{} sepanjang $\gamma$ dan

\mavergleichskette
{\vergleichskette
{ a
}
{ \in }{ \R
}
{  }{
}
{    }{
}
{   }{
}
}
{}{}{,}
maka $aF$ juga merupakan medan vektor paralel.
} {Jika $F$ dan $G$ adalah medan-medan vektor paralel sepanjang $\gamma$, maka
\mathl{F +G}{} juga merupakan medan vektor paralel.
}}
\faktzusatz {}
\faktzusatz {}

 }
{ Lihat [[Hyperfläche/Kurve/Paralleles Vektorfeld/Eigenschaften/Fakt/Beweis/Aufgabe|Soal 6.3]]. }

Untuk suatu kurva terdiferensialkan
\maabbdisp {\gamma} { I } { Y
} {}
medan kecepatan $\gamma'$ khususnya merupakan suatu medan vektor sepanjang $\gamma$; karena itu, konsep apakah medan ini paralel dapat diterapkan.

__NOEDITSECTION__

\inputfaktbeweis
{Differenzierbare Hyperfläche/Differenzierbare Kurve/Ableitung parallel/Geodätische/Fakt}
{Lemma}
{}
{

\faktsituation {Misalkan

\mavergleichskette
{\vergleichskette
{ Y
}
{ \subseteq }{ W
}
{ \subseteq }{ \R^n
}
{    }{
}
{   }{
}
}
{}{}{,}
$W$
\definitionsverweis {terbuka}{}{,}
dan $Y$ suatu
\definitionsverweis {hipermuka terdiferensialkan}{}{}
dan misalkan
\maabb {\gamma} { I } { Y
} {}
suatu kurva yang
\definitionsverweis {dua kali terdiferensialkan secara kontinu}{}{.}}
\faktfolgerung {Maka $\gamma$ merupakan suatu
\definitionsverweis {kurva geodesik}{}{}
tepat ketika medan turunannya
\zusatzklammer {medan kecepatannya} {} {} $\gamma'$ merupakan suatu
\definitionsverweis {medan vektor paralel}{}{}
sepanjang $\gamma$.}
\faktzusatz {}
\faktzusatz {}

}
{

Misalkan
\maabbdisp {\gamma^{\prime \prime}} { I } { \R^n
} {}
adalah turunan kedua. Berdasarkan definisi, $\gamma$ merupakan suatu kurva geodesik jika
\mathl{\gamma^{\prime \prime} (t)}{} selalu tegak lurus terhadap ruang singgung $T_{\gamma(t)} Y$, yang berlaku tepat ketika proyeksi ortogonal $\gamma^{\prime \prime}(t)$ pada ruang singgung sama dengan $0$. Hal ini ekuivalen dengan

\mavergleichskettedisp
{\vergleichskette
{ (\nabla_\gamma \gamma')(t)
}
{ =} { \pi_{\gamma(t)} \left(  \gamma^{\prime \prime} (t)  \right)
}
{ =} { 0
}
{ } {
}
{ } {
}
}
{}{}{}
untuk semua

\mavergleichskette
{\vergleichskette
{ t
}
{ \in }{ I
}
{  }{
}
{    }{
}
{   }{
}
}
{}{}{,}
yang berarti bahwa $\gamma'$ merupakan suatu
\definitionsverweis {medan vektor paralel}{}{}
sepanjang $\gamma$.

}

Pernyataan berikut menjadi dasar bagi keberadaan transpor paralel.
__NOEDITSECTION__

\inputfaktbeweis
{Hyperfläche/Kurve/Anfangstangentenvektor/Paralleles Vektorfeld/Existenz/Fakt}
{Teorema}
{}
{

\faktsituation {Misalkan

\mavergleichskette
{\vergleichskette
{ Y
}
{ \subseteq }{ W
}
{ \subseteq }{\R^n
}
{    }{
}
{   }{
}
}
{}{}{,}
$W$
\definitionsverweis {terbuka}{}{,}
dan $Y$ suatu
\definitionsverweis {hipermuka terdiferensialkan}{}{} kelas $C^2$
dan misalkan
\maabb {\gamma} { I } { Y
} {}
suatu
\definitionsverweis {kurva terdiferensialkan}{}{.}
Misalkan

\mavergleichskette
{\vergleichskette
{ t_0
}
{ \in }{ I
}
{  }{
}
{    }{
}
{   }{
}
}
{}{}{,}

\mavergleichskette
{\vergleichskette
{ \gamma(t_0)
}
{ = }{ P
}
{  }{
}
{    }{
}
{   }{
}
}
{}{}{}
dan misalkan

\mavergleichskette
{\vergleichskette
{ v
}
{ \in }{ T_PY
}
{  }{
}
{    }{
}
{   }{
}
}
{}{}{}
suatu
\definitionsverweis {vektor singgung}{}{.}}
\faktfolgerung {Maka terdapat suatu
\definitionsverweis {medan vektor tangensial}{}{}
$F$ yang ditentukan secara unik sepanjang $\gamma$, yang bersifat
\definitionsverweis {paralel}{}{}
dan memenuhi

\mavergleichskettedisp
{\vergleichskette
{ F(t_0)
}
{ =} { v
}
{ } {
}
{ } {
}
{ } {
}
}
{}{}{.}
}
\faktzusatz {}
\faktzusatz {}

}
{

Pilih suatu medan normal satuan terdiferensialkan $N$ sepanjang $\gamma$; regularitas $C^2$ pada $Y$ menjamin pilihan lokal tersebut dapat diteruskan sepanjang interval $I$. Kita mencari suatu pemetaan terdiferensialkan
\maabbdisp {F} { I } { \R^n
} {,}
yang memenuhi ketiga syarat
\aufzaehlungdrei{
\mavergleichskettedisp
{\vergleichskette
{ F(t)
}
{ \in} { T_{\gamma(t)} Y
}
{ } {
}
{ } {
}
{ } {
}
}
{}{}{}
untuk semua

\mavergleichskette
{\vergleichskette
{ t
}
{ \in }{ I
}
{  }{
}
{    }{
}
{   }{
}
}
{}{}{,}
}{
\mavergleichskettedisp
{\vergleichskette
{ F'(t)
}
{ \in} { N_{\gamma(t)} Y
}
{ } {
}
{ } {
}
{ } {
}
}
{}{}{}
untuk semua

\mavergleichskette
{\vergleichskette
{ t
}
{ \in }{ I
}
{  }{
}
{    }{
}
{   }{
}
}
{}{}{,}
}{
\mavergleichskettedisp
{\vergleichskette
{ F (t_0)
}
{ =} { v
}
{ } {
}
{ } {
}
{ } {
}
}
{}{}{.}
}
Syarat pertama berarti bahwa $F(t)$ tegak lurus terhadap vektor normal satuan $N (\gamma(t))$, sehingga

\mavergleichskettedisp
{\vergleichskette
{ \left\langle F( t) ,  N( \gamma(t))   \right\rangle
}
{ =} { 0
}
{ } {
}
{ } {
}
{ } {
}
}
{}{}{.}
Oleh karena itu, berlaku pula

\mavergleichskettedisp
{\vergleichskette
{ 0
}
{ =} { \left\langle F(t)  ,  N( \gamma(t))   \right\rangle'
}
{ =} { \left\langle F'(t)  ,  N( \gamma(t))   \right\rangle +  \left\langle F(t)  ,  (N( \gamma(t)))'   \right\rangle
}
{ } {
}
{ } {
}
}
{}{}{.}
Syarat kedua, bahwa $F'(t)$ berada di ruang normal, berarti

\mavergleichskettedisp
{\vergleichskette
{ F'(t) - \left\langle F'( t) ,  N( \gamma(t))   \right\rangle N(\gamma(t) )
}
{ =} { 0
}
{ } {
}
{ } {
}
{ } {
}
}
{}{}{}
untuk semua

\mavergleichskette
{\vergleichskette
{ t
}
{ \in }{ I
}
{  }{
}
{    }{
}
{   }{
}
}
{}{}{,}
yang dengan bantuan persamaan sebelumnya dapat kita ubah menjadi

\mavergleichskettedisp
{\vergleichskette
{ F'(t) + \left\langle F( t) ,  (N( \gamma(t)))'   \right\rangle N(\gamma(t) )
}
{ =} { 0
}
{ } {
}
{ } {
}
{ } {
}
}
{}{}{}
untuk semua

\mavergleichskette
{\vergleichskette
{ t
}
{ \in }{ I
}
{  }{
}
{    }{
}
{   }{
}
}
{}{}{}
atau menjadi

\mavergleichskettedisp
{\vergleichskette
{ F'(t)
}
{ =} { - \left\langle F( t) ,  (N( \gamma(t)))'   \right\rangle N(\gamma(t) )
}
{ } {
}
{ } {
}
{ } {
}
}
{}{}{}
untuk semua

\mavergleichskette
{\vergleichskette
{ t
}
{ \in }{ I
}
{  }{
}
{    }{
}
{   }{
}
}
{}{}{.}
Di sini terdapat suatu persamaan diferensial linear biasa orde pertama untuk $F$ yang, bersama dengan syarat awal

\mavergleichskette
{\vergleichskette
{ F(t_0)
}
{ = }{ v
}
{  }{
}
{    }{
}
{   }{
}
}
{}{}{,}
menurut
[[Picard Lindelöf/Lokal Lipschitz/Lokale Existenz und Eindeutigkeit/Fakt|Teorema 56.2 (Analisis (Osnabrück 2021--2023))]],
memiliki solusi lokal tunggal. Karena persamaan ini linear dengan koefisien kontinu pada $I$, solusi tersebut dapat dilanjutkan secara unik ke seluruh interval $I$. Karena $v\in T_PY$, nilai awal memenuhi $\langle v,N(P)\rangle=0$. Sepanjang solusi persamaan diferensial tersebut,
\[
\frac{d}{dt}\left\langle F(t),N(\gamma(t))\right\rangle
=-\left\langle F(t),(N\circ\gamma)'(t)\right\rangle
+\left\langle F(t),(N\circ\gamma)'(t)\right\rangle=0.
\]
Jadi $F(t)$ tetap tangensial dan persamaan diferensial membuat $F'(t)$ tetap normal. Dengan demikian, solusi tersebut memang merupakan medan paralel yang diminta dan keunikan persamaan diferensial memberi keunikannya.

\emph{Catatan edisi: Sumber hanya mengasumsikan hipermuka terdiferensialkan, menerapkan teorema eksistensi lokal langsung pada seluruh interval $I$, dan tidak membuktikan bahwa solusi persamaan diferensial tetap tangensial. Hipotesis diperkuat menjadi $C^2$, kelanjutan global untuk sistem linear dinyatakan, medan normal diperkenalkan secara eksplisit, dan perhitungan invarian yang hilang dilengkapi.}

}

\inputbemerkung
{}
{

Dalam situasi
[[Hyperfläche/Kurve/Anfangstangentenvektor/Paralleles Vektorfeld/Existenz/Fakt|Teorema 6.9]],
persamaan diferensial yang harus diselesaikan adalah sistem

\mavergleichskettedisp
{\vergleichskette
{ F_i'(t)
}
{ =} { - \left(  \sum_{j = 1}^n  F_j (t) (N ( \gamma(t)))'_j  \right)   N(\gamma(t))_i
}
{ } {
}
{ } {
}
{ } {
}
}
{}{}{}
untuk

\mavergleichskette
{\vergleichskette
{ i
}
{ = }{ 1 , \ldots , n
}
{  }{
}
{    }{
}
{   }{
}
}
{}{}{}
atau, dalam notasi matriks,

\mavergleichskettedisp
{\vergleichskette
{ \begin{pmatrix} F_1' \\ \vdots\\ F_n'   \end{pmatrix}
}
{ =} { - A \begin{pmatrix} F_1  \\ \vdots\\ F_n   \end{pmatrix}
}
{ } {
}
{ } {
}
{ } {
}
}
{}{}{}
dengan entri-entri

\mavergleichskettedisp
{\vergleichskette
{ A_{ij} (t)
}
{ =} { N(\gamma(t))_i (N(\gamma(t)))'_j
}
{ } {
}
{ } {
}
{ } {
}
}
{}{}{.}

\emph{Catatan edisi: Satu tanda kurung tutup berlebih pada faktor $(N(\gamma(t)))'_j$ dihapus, sehingga koefisiennya adalah $A_{ij}=N_iN'_j$.}

}

\inputbeispiel{}
{

Kita meninjau seperempat lingkaran besar
\maabbeledisp {} {[0 , \pi/2] } { \R^3
} { t } { \left(  \cos  t   , \,   \sin  t  , \,   0                 \right)
} {,}
pada permukaan bola satuan $S$. Berlaku

\mavergleichskette
{\vergleichskette
{ \gamma(0)
}
{ = }{ \left(  1   , \,   0  , \,   0                 \right)
}
{ = }{ P
}
{    }{
}
{   }{
}
}
{}{}{}
dan

\mavergleichskette
{\vergleichskette
{ \gamma(\pi/2)
}
{ = }{ \left(  0   , \,   1  , \,   0                 \right)
}
{ = }{ Q
}
{    }{
}
{   }{
}
}
{}{}{}
dengan
\definitionsverweis {ruang-ruang singgung}{}{}

\mavergleichskette
{\vergleichskette
{ T_PS
}
{ = }{ \R e_2 +\R e_3
}
{  }{
}
{    }{
}
{   }{
}
}
{}{}{}
dan

\mavergleichskette
{\vergleichskette
{ T_QS
}
{ = }{ \R e_1 +\R e_3
}
{  }{
}
{    }{
}
{   }{
}
}
{}{}{.}
Vektor normal satuan yang mengarah ke dalam sepanjang lintasan tersebut adalah

\mavergleichskette
{\vergleichskette
{ N(t)
}
{ = }{ -\gamma(t)
}
{  }{
}
{    }{
}
{   }{
}
}
{}{}{.}
Matriks dari
[[Hyperfläche/Kurve/Anfangstangentenvektor/Paralleles Vektorfeld/Differentialgleichung/Explizite Gestalt/Bemerkung|Catatan 6.10]],
yang mendeskripsikan persamaan diferensial bagi suatu medan vektor paralel, adalah

\mathdisp {\begin{pmatrix}  -\sin  t \cos  t   &   \cos  t \cos  t  &  0  \\  - \sin  t  \sin  t  &  \sin  t  \cos  t  &  0  \\ 0 &  0  &  0 \end{pmatrix}} { . }

Fungsi konstan

\mavergleichskettedisp
{\vergleichskette
{ F(t)
}
{ =} { e_3
}
{ =} { \begin{pmatrix}  0  \\ 0\\  1   \end{pmatrix}
}
{ } {
}
{ } {
}
}
{}{}{}
merupakan suatu solusi. Selanjutnya,

\mavergleichskettedisp
{\vergleichskette
{ G(t)
}
{ =} { \begin{pmatrix}  - \sin  t  \\ \cos  t \\  0   \end{pmatrix}
}
{ } {
}
{ } {
}
{ } {
}
}
{}{}{}
merupakan suatu solusi; di satu pihak,

\mavergleichskettedisp
{\vergleichskette
{ G'(t)
}
{ =} { \begin{pmatrix}  - \cos  t  \\ - \sin  t \\  0   \end{pmatrix}
}
{ } {
}
{ } {
}
{ } {
}
}
{}{}{}
dan, di pihak lain,

\mavergleichskettedisphandlinks
{\vergleichskettedisphandlinks
{ - \begin{pmatrix}  -\sin  t \cos  t   &   \cos  t \cos  t  &  0  \\  - \sin  t  \sin  t  &  \sin  t  \cos  t  &  0  \\ 0 &  0  &  0 \end{pmatrix} \begin{pmatrix}  - \sin  t  \\ \cos  t \\  0   \end{pmatrix}
}
{ =} { - \begin{pmatrix}  \sin^{ 2  }  t \cos  t + \cos^{ 3  }  t  \\ \sin^{ 3  }  t +  \sin  t \cos^{ 2  }  t \\  0   \end{pmatrix}
}
{ =} { - \begin{pmatrix}  \cos  t  \\ \sin  t \\  0   \end{pmatrix}
}
{ } {
}
{ } {
}
}
{}{}{.}

}

  \zwischenueberschrift{Transpor Paralel}

Misalkan

\mavergleichskette
{\vergleichskette
{ Y
}
{ \subseteq }{ W
}
{ \subseteq }{ \R^n
}
{    }{
}
{   }{
}
}
{}{}{,}
$W$
\definitionsverweis {terbuka}{}{,}
dan $Y$ suatu
\definitionsverweis {hipermuka terdiferensialkan}{}{} kelas $C^2$.
Untuk titik-titik

\mavergleichskette
{\vergleichskette
{ P,Q
}
{ \in }{ Y
}
{  }{
}
{    }{
}
{   }{
}
}
{}{}{}
tidak terdapat hubungan alami antara
\definitionsverweis {ruang singgung}{}{}
$T_PY$ dan ruang singgung $T_QY$. Misalkan
\maabbdisp {\gamma} {[a,b]} {Y
} {}
suatu
\definitionsverweis {kurva terdiferensialkan}{}{}
di $Y$ dengan

\mavergleichskette
{\vergleichskette
{ \gamma(a)
}
{ = }{ P
}
{  }{
}
{    }{
}
{   }{
}
}
{}{}{}
dan

\mavergleichskette
{\vergleichskette
{ \gamma(b)
}
{ = }{ Q
}
{  }{
}
{    }{
}
{   }{
}
}
{}{}{.}
Kita dapat bertanya apakah sepanjang lintasan ini terdapat hubungan yang bermakna antara ruang-ruang singgung tersebut.

Untuk suatu kurva terdiferensialkan tetap $\gamma$ di $Y$ yang menghubungkan titik-titik
\mathkor {} {P =\gamma(a)} {dan} {Q = \gamma(b)} {,}
berdasarkan
[[Hyperfläche/Kurve/Anfangstangentenvektor/Paralleles Vektorfeld/Existenz/Fakt|Teorema 6.9]],
ditentukan suatu pemetaan
\maabbdisp {} {T_PY} {T_QY
} {.}
Pemetaan ini memetakan vektor awal

\mavergleichskette
{\vergleichskette
{ v
}
{ \in }{ T_PY
}
{  }{
}
{    }{
}
{   }{
}
}
{}{}{}
ke vektor

\mavergleichskette
{\vergleichskette
{ F(b)
}
{ \in }{ T_QY
}
{  }{
}
{    }{
}
{   }{
}
}
{}{}{,}
dengan $F$ adalah medan vektor paralel yang ditentukan secara unik sepanjang $\gamma$. Pemetaan ini disebut \stichwort {transpor paralel} {} sepanjang $\gamma$. Pemetaan tersebut dilambangkan dengan
\maabbdisp {\Psi_\gamma} { T_{\gamma(a)}Y} {T_{\gamma(b)}Y
} {.}

\bild{ \begin{center}
\includegraphics[width=5.5cm]{ \bildeinlesungsvg {Parallel transport sphere2} {svg} }
\end{center}
\bildtext {} }

\bildlizenz { Parallel transport sphere2.svg } {} {Silly rabbit} {en Wikipedia} {CC-by-sa 3.0} {}

__NOEDITSECTION__

\inputfaktbeweis
{Hyperfläche/Kurve/Paralleltransport/Isometrie/Fakt}
{Teorema}
{}
{

\faktsituation {Misalkan

\mavergleichskette
{\vergleichskette
{ Y
}
{ \subseteq }{ W
}
{ \subseteq }{ \R^n
}
{    }{
}
{   }{
}
}
{}{}{,}
$W$
\definitionsverweis {terbuka}{}{,}
dan $Y$ suatu
\definitionsverweis {hipermuka terdiferensialkan}{}{} kelas $C^2$
dan misalkan
\maabb {\gamma} {[a,b]} {Y
} {}
suatu
\definitionsverweis {kurva terdiferensialkan}{}{}
dengan

\mavergleichskette
{\vergleichskette
{ \gamma(a)
}
{ = }{ P
}
{  }{
}
{    }{
}
{   }{
}
}
{}{}{}
dan

\mavergleichskette
{\vergleichskette
{ \gamma(b)
}
{ = }{ Q
}
{  }{
}
{    }{
}
{   }{
}
}
{}{}{.}}
\faktfolgerung {Maka
\definitionsverweis {transpor paralel}{}{}
sepanjang $\gamma$ merupakan suatu
\definitionsverweis {isometri}{}{}
\maabbdisp {} {T_PY} {T_QY
} {.}}
\faktzusatz {}
\faktzusatz {}

}
{

Untuk membuktikan linearitas, misalkan

\mavergleichskette
{\vergleichskette
{ v,w
}
{ \in }{ T_PY
}
{  }{
}
{    }{
}
{   }{
}
}
{}{}{}
dan

\mavergleichskette
{\vergleichskette
{ r,s
}
{ \in }{ \R
}
{  }{
}
{    }{
}
{   }{
}
}
{}{}{}
diberikan. Misalkan
\mathkor {} {F} {dan} {G} {}
adalah medan-medan vektor paralel sepanjang $\gamma$ yang ditentukan secara unik menurut
[[Hyperfläche/Kurve/Anfangstangentenvektor/Paralleles Vektorfeld/Existenz/Fakt|Teorema 6.9]]
dengan

\mavergleichskette
{\vergleichskette
{ F(a)
}
{ = }{ v
}
{  }{
}
{    }{
}
{   }{
}
}
{}{}{}
dan

\mavergleichskette
{\vergleichskette
{ G(a)
}
{ = }{ w
}
{  }{
}
{    }{
}
{   }{
}
}
{}{}{.}
Menurut
[[Hyperfläche/Kurve/Paralleles Vektorfeld/Eigenschaften/Fakt|Lemma 6.7]],
\mathl{rF+sG}{} merupakan suatu medan vektor paralel dengan

\mavergleichskette
{\vergleichskette
{ (rF+sG)(a)
}
{ = }{ rv+sw
}
{  }{
}
{    }{
}
{   }{
}
}
{}{}{.}
\emph{Catatan edisi: Nilai awal kombinasi linear diperbaiki dari $v+w$ menjadi $rv+sw$.}
Berdasarkan keunikan dari
[[Hyperfläche/Kurve/Anfangstangentenvektor/Paralleles Vektorfeld/Existenz/Fakt|Teorema 6.9]],
\mathl{rF+sG}{} dengan demikian merupakan medan vektor paralel untuk vektor singgung
\mathl{rv+sw}{.} Oleh karena itu,

\mavergleichskettedisp
{\vergleichskette
{ \Psi_{\gamma} (rv+sw)
}
{ =} { (rF+sG)(b)
}
{ =} { r F(b) +s G(b)
}
{ =} { r \Psi_{\gamma} (v) + s \Psi_{\gamma} (w)
}
{ } {
}
}
{}{}{.}

Untuk membuktikan kompatibilitas dengan hasil kali dalam, misalkan kembali

\mavergleichskette
{\vergleichskette
{ v,w
}
{ \in }{ T_PY
}
{  }{
}
{    }{
}
{   }{
}
}
{}{}{}
diberikan, dan misalkan $F,G$ adalah medan-medan vektor paralel yang bersesuaian. Berlaku

\mavergleichskettedisp
{\vergleichskette
{ \left\langle  F(t) , G(t)  \right\rangle'
}
{ =} { \left\langle  F'(t) , G(t)  \right\rangle  +  \left\langle  F(t) , G'(t)  \right\rangle
}
{ =} { 0
}
{ } {
}
{ } {
}
}
{}{}{,}
karena $F,G$ bersifat tangensial, sedangkan $F',G'$ ortogonal terhadap ruang singgung. Jadi,
\mathl{\left\langle  F(t) , G(t)  \right\rangle}{} konstan sepanjang lintasan. Oleh karena itu,

\mavergleichskettedisp
{\vergleichskette
{ \left\langle \Psi_\gamma(v)  , \Psi_\gamma(w)  \right\rangle
}
{ =} { \left\langle  F(b)  ,  G(b)   \right\rangle
}
{ =} { \left\langle  F(a)  ,  G(a)   \right\rangle
}
{ =} { \left\langle  v  , w  \right\rangle
}
{ } {
}
}
{}{}{.}

\emph{Catatan edisi: Faktor kedua pada hasil kali dalam di titik awal diperbaiki dari $F(a)$ menjadi $G(a)$.}
Bijektivitasnya juga jelas dari sini.

}

\inputbeispiel{}
{

Kita meninjau seperempat lingkaran besar
\maabbeledisp {} {[0 , \pi/2] } { \R^3
} { t } { \left(  \cos  t   , \,   \sin  t  , \,   0                 \right)
} {,}
pada permukaan bola satuan $S$ dan melanjutkan pembahasan dari
[[Kugeloberfläche/Viertelgroßkreis/Parallele Vektorfelder/Beispiel|Contoh 6.11]].
Berlaku

\mavergleichskette
{\vergleichskette
{ \gamma(0)
}
{ = }{ \left(  1   , \,   0  , \,   0                 \right)
}
{ = }{ P
}
{    }{
}
{   }{
}
}
{}{}{}
dan

\mavergleichskette
{\vergleichskette
{ \gamma(\pi/2)
}
{ = }{ \left(  0   , \,   1  , \,   0                 \right)
}
{ = }{ Q
}
{    }{
}
{   }{
}
}
{}{}{}
dengan
\definitionsverweis {ruang-ruang singgung}{}{}

\mavergleichskette
{\vergleichskette
{ T_PS
}
{ = }{ \R e_2 +\R e_3
}
{  }{
}
{    }{
}
{   }{
}
}
{}{}{}
dan

\mavergleichskette
{\vergleichskette
{ T_QS
}
{ = }{ \R e_1 +\R e_3
}
{  }{
}
{    }{
}
{   }{
}
}
{}{}{.}
Berdasarkan perhitungan dalam contoh yang dirujuk tersebut, untuk
\definitionsverweis {transpor paralel}{}{}
berlaku

\mavergleichskette
{\vergleichskette
{ \Psi_\gamma(e_3)
}
{ = }{ e_3
}
{  }{
}
{    }{
}
{   }{
}
}
{}{}{}
dan

\mavergleichskette
{\vergleichskette
{ \Psi_\gamma(e_2)
}
{ = }{ - e_1
}
{  }{
}
{    }{
}
{   }{
}
}
{}{}{.}

}

Jika
\maabbdisp {\gamma} { I } { Y
} {}
merupakan suatu kurva kontinu yang terdiferensialkan sepotong-sepotong, transpor paralel sepanjang $\gamma$ didefinisikan dengan melaksanakan secara berurutan transpor-transpor paralel pada setiap potongan kurva yang terdiferensialkan.

 [[Kategorie:Latexseite]]
